\chapter{图像识别与检索}
\label{chap:vision-image-recognition}

一张路口照片中有两辆车。给出“街道”这个场景名，尚未说明车辆在哪里；给出两个框，尚未说明哪辆是红色、是否存在遮挡、距离相机多远；找到另一张相似街景，也未必找到了同一辆车。本章围绕这些逐渐具体的请求，建立从像素到语义、范围、测量和对应关系的计算。

通用的卷积、注意力和表示学习已经提供了特征提取工具。在这里，需要进一步说明特征对应整幅图、一个对象还是一个局部位置，监督约束哪个输出，以及模型预测如何还原成用户需要的结果。街景贯穿对象与空间任务，正常零件与划痕则用于说明没有完整缺陷类别表时怎样判断异常。以下算例均为教学构造，数值用于复核机制。

\section{语义与属性判定}
\label{sec:vision-recognition-semantics}

类别说明“属于什么”，属性说明“具有什么特征”，异常判断说明“是否偏离规定的正常状态”。三种输出可以来自同一表示，但需要不同的标签和判定依据。场景分类把整幅画面的环境概括为街道或厨房，对象分类则需要明确被判断的对象；厨房这一场景名不能展开为一份完整的物品清单。

\subsection{对象类别与场景判定}
\label{sec:vision-recognition-classification}

在单标签任务中，一个样本只对应类别表中的一项；多标签任务允许多项同时成立。例如车辆裁剪图可按主要车型给一个标签，整张街景却可以同时含公交车、轿车和行人。损失函数必须遵守这一区别：互斥类别可用softmax，多种可共存对象通常分别输出是否存在的得分。

\subsubsection{已有视觉表示与分类头：DINOv3的任务接入}
\label{sec:vision-recognition-dinov3}

已有编码器可以把图像转成全局向量及保留位置的特征。DINOv3提供了这类可复用表示；本节关心的是拿到表示后怎样完成有标签的任务，其完整自监督训练另属表示学习。全局向量压缩整图信息，空间特征仍与图像块位置相连，二者适合接入不同的预测头。\nocite{vision-2508-10104}\href{https://arxiv.org/abs/2508.10104}{DINOv3技术报告}给出了其表示与下游任务设置。

设冻结编码器$f_{\bm{\theta}}$得到$\bm{z}=f_{\bm{\theta}}(\bm{X})\in\mathbb{R}^{d}$，分类头包含$\bm{W}\in\mathbb{R}^{k\times d}$与$\bm{b}\in\mathbb{R}^{k}$。对$k$个互斥类别，计算
\begin{equation}
 p(y=r\mid\bm{X})=
 \frac{\exp(\bm{w}_r^\top\bm{z}+b_r)}
 {\sum_{j=1}^{k}\exp(\bm{w}_j^\top\bm{z}+b_j)},
 \qquad
 L=-\frac1n\sum_{i=1}^{n}\log p(y_i\mid\bm{X}_i)\text{。}
 \label{eq:vision-linear-head}
\end{equation}
这里$\bm{w}_r^\top$是$\bm{W}$的第$r$行。冻结训练只更新$\bm{W},\bm{b}$；全量适配还更新$\bm{\theta}$，所需梯度存储和训练数据条件随之变化。使用时保留编码器、分类头以及相同的预处理，取最大得分类别，或按事先验证的规则拒绝不确定输入。

假设三类车型的得分为$(2,1,0)$。经过softmax计算，概率依次约为0.665、0.245和0.090。它们是在给定类别集合内归一化的结果，不能证明图像中一定存在可辨车辆。若编码器读取整幅街景，得分还可能依赖道路背景；若任务要判定左车的车型，应先取得左车区域，或使用明确的区域监督。

外观接近的物种分类揭示了另一种困难：全局背景足以区分拍摄地点，却可能不足以区分鸟喙或叶缘。iNaturalist包含细粒度类别和不均衡样本，评价时应同时检查小部位分辨率和类别覆盖\scite{vision-1707-06642}
例如常见类别有90张且全对，少见类别有10张且全错，样本平均准确率为90\%，两个类别的平均准确率却为50\%。这两项分别强调观测频率与类别公平覆盖，不能互相替代。

\subsubsection{文字指定类别与零样本识别：CLIP接口}
\label{sec:vision-recognition-clip}

如果使用时才给出“轿车、货车、公交车”，固定输出维度的分类头不能直接容纳新类别。图文表示把类别描述也编码为向量，使图像可以与使用时给定的候选比较。CLIP的双编码器训练将在\ref{sec:vision-models-contrastive}节展开；这里先把两侧编码器视为已训练的接口。

令$\bar{\bm{z}}$与$\bar{\bm{t}}_r$分别为单位长度的图像向量和第$r$条描述向量，按$s_r=\bar{\bm{z}}^\top\bar{\bm{t}}_r$排序。使用温度$\tau>0$可计算候选内归一化得分
\begin{equation}
 q_r=\frac{\exp(s_r/\tau)}
 {\sum_{j\in C}\exp(s_j/\tau)}\text{，}
 \label{eq:vision-clip-classification}
\end{equation}
其中$C$是这次输入的候选索引集合。分母随候选变化，$q_r$不是不依赖候选的存在概率。可以对“这是一张公交车照片”等多个模板的单位文字向量求平均后再归一化，但模板集合须在评价前固定。

设四项相似度为$0.31,0.28,0.19,0.12$，最后一项是不在图中的“帆船”。最大值只使第一项在这四项中领先。把前三项全部换成不相关词，模型仍会返回一个最大值。开放使用要另设含未知对象的验证协议、拒绝阈值或目标存在判断；单凭相似度最大不能完成拒识。整图相似也无法保证“红色”绑定到货车而非旁边的轿车。

\subsection{对象属性判定}
\label{sec:vision-recognition-attributes}

要记录“红色货车”和“蓝色轿车”，仅输出整图属性集合$\{\text{红色},\text{蓝色}\}$会丢掉归属。可以先定位对象，再对各区域提取表示，形成$(\text{对象编号},\text{属性},\text{取值})$记录。位置与属性共享编号，使后来计数、搜索或人工核验仍指向同一个对象。

属性之间是否互斥应由标签定义决定。“红色车身”和“车灯开启”可以同时成立，“完整车身的主色”则可以规定为一个单标签字段。对于可共存的$k$个二值属性，令$p_{i,r}$为预测，$a_{i,r}\in\{0,1\}$为标注，$m_{i,r}$表示该属性是否有可靠标注，则可用
\begin{equation}
 L_{\mathrm{attr}}=
 -\frac{\sum_{i,r}m_{i,r}
 [a_{i,r}\log p_{i,r}+(1-a_{i,r})\log(1-p_{i,r})]}
 {\sum_{i,r}m_{i,r}}\text{。}
 \label{eq:vision-attribute-loss}
\end{equation}
分母要求至少有一个有效标注；全无标注的批次不计算此项。把“未标注”直接置为“不存在”，会训练出系统性的假负例。

裁剪区域过小会切掉车灯，过大又可能纳入旁车颜色。应检查属性所需部位是否可见，并在不可判定时保留未知状态。完整对象编号并不能自行修正错误的视觉证据，但它使属性错绑能够被单独定位。

\subsection{正常状态与异常判定}
\label{sec:vision-recognition-anomaly}

工业零件的异常类型可能不断出现，训练集却主要含正常样本。此时可学习或保存正常外观，再度量新观测与正常外观的偏离。重建误差、分布密度和局部近邻距离是不同的偏离依据，得分都需要结合独立验证集解释。

\subsubsection{正常局部特征记忆与近邻评分：PatchCore}
\label{sec:vision-recognition-patchcore}

PatchCore从已有编码器的中间层提取并局部聚合特征，将正常图像中各位置的向量存入记忆。相邻尺度的特征对齐后保留局部上下文，使用时不需要先列出所有划痕和破损类别\scite{vision-2106-08265}

全部正常片段可能很大，因而从特征集合$\bm{M}$中保留代表子集$\bm{M}_{\mathrm{c}}$。其覆盖目标可写为
\begin{equation}
 \min_{\substack{\bm{M}_{\mathrm{c}}\subseteq\bm{M}\\
 |\bm{M}_{\mathrm{c}}|=m}}
 \max_{\bm{z}\in\bm{M}}\min_{\bm{u}\in\bm{M}_{\mathrm{c}}}
 \lVert\bm{z}-\bm{u}\rVert_2\text{。}
 \label{eq:vision-patchcore-coreset}
\end{equation}
原方法使用近似贪心选择，而非精确穷举全部子集。已选一个起点后，反复加入离当前子集最远的特征，可以保留正常模式中的不同部分；这个过程在建立记忆，不是在用缺陷标签重新训练编码器。

查询图像位置$p$的基本局部分数为$d_p=\min_{\bm{u}\in\bm{M}_{\mathrm{c}}}\lVert\bm{z}_p-\bm{u}\rVert_2$。将这些分数按空间排列、插值和平滑可形成异常热图。最大局部距离$s_*=\max_p d_p$是整图评分的起点；原方法还根据最异常片段$\bm{z}_*$所对应正常近邻$\bm{u}_*$周围的支持情况重新加权。令$\bm{N}_b(\bm{u}_*)$包含$\bm{u}_*$及其在记忆中的近邻，则
\begin{equation}
 s=\left(1-
 \frac{\exp\lVert\bm{z}_*-\bm{u}_*\rVert_2}
 {\sum_{\bm{u}\in\bm{N}_b(\bm{u}_*)}
  \exp\lVert\bm{z}_*-\bm{u}\rVert_2}\right)s_*\text{。}
 \label{eq:vision-patchcore-score}
\end{equation}
邻域大小$b$是需固定的设置，实际应大于1；否则分子分母相同，分数退化为零。若最近正常片段较孤立，其他邻居的距离更大，括号中的系数更接近1，最大距离受到的压低更少。这个重加权改变整图分数，不改变局部距离图的来源。

假设一维正常记忆为$\{0,1,4\}$，查询局部特征为$\{0.2,0.9,6\}$，对应距离是$0.2,0.1,2$，最大值由第三个片段给出。若$b=2$，取最近正常片段4及其邻居1，整图分数为$2[1-\exp(2)/(\exp(2)+\exp(5))]\approx1.905$。区域线索保留了小范围的偏离，但若新设备使所有特征整体平移，正常零件也可能远离记忆。高分只能表示对当前正常参照的偏离，不能直接命名缺陷或判定业务不合格。

\section{对象定位与区域划分}
\label{sec:vision-recognition-regions}

定位要求把判断接回图上的位置。边界框给出大致范围，掩码给出像素归属，透明或混合边界还需要连续贡献比例。希望处理哪个目标，可以由固定类别、文字或点框提示指定；目标选择与范围形成需要分别解释。

\subsection{对象范围的定位}
\label{sec:vision-recognition-detection}

\term{边界框}（bounding box）用一个矩形包围目标。以下采用连续边界$\bm{b}=(x_1,y_1,x_2,y_2)$，宽高分别为$x_2-x_1$和$y_2-y_1$，不加离散像素计数中的1。检测结果是多个$(\bm{b},c,s)$组成的集合，其中$c$是类别，$s$是得分。标签应说明框包围可见部分还是完整对象；旋转目标也可改用中心、宽高和角度表示。

比较两个框的重叠程度时，\term{交并比}（intersection over union, IoU）将交集面积除以并集面积：
\begin{equation}
 \operatorname{IoU}(\bm{b},\bm{b}')
 =\frac{|\bm{b}\cap\bm{b}'|}
 {|\bm{b}|+|\bm{b}'|-|\bm{b}\cap\bm{b}'|}\text{。}
 \label{eq:vision-iou}
\end{equation}
这里矩形应有正面积，无交集时分子为零。取两个面积均为800的框，边界分别为$(10,10,50,30)$和$(20,10,60,30)$。它们的交集面积为600，IoU为$600/1000=0.6$。若它们是对同一辆车的两个预测，应去重；若属于两辆互相遮挡的车，直接去掉一个则会漏检。

\subsubsection{整图网格预测：YOLO}
\label{sec:vision-recognition-yolo}

原始YOLO把图像划为$S\times S$网格，每个单元预测$B$个框和一组共享的条件类别概率。目标中心落在哪个单元，就由该单元承担监督。框中心相对于单元定位，宽高相对于整图归一化；每框置信度学习目标存在与框质量的联合信号\scite{vision-1506-02640}

若单元列行为$(j,i)$，预测中心偏移为$(u,v)$，输入宽高为$(w,h)$，则中心还原为$((j+u)w/S,(i+v)h/S)$。训练在负责单元的多个预测框中选择当前IoU最高者承担该对象的坐标监督。原始平方误差目标分别约束中心、宽高的平方根、置信度和类别；没有负责对象的候选也接受置信度监督，但权重较低，以免大量背景压倒对象信号。

每框置信度写为$\mathbb{P}(\mathrm{obj})\operatorname{IoU}$，乘上共享的$\mathbb{P}(c\mid\mathrm{obj})$得到类别相关得分。它将语义与定位质量合在一起，不能只解释为类别概率。使用时解码全部框，按阈值过滤，再用\term{非极大值抑制}（non-maximum suppression, NMS）去掉同类中与当前高分框过度重叠的低分框。

NMS按得分从高到低取一个框，保留它并删除重叠超过阈值的候选，再对余下框重复，直到没有候选。对于前述IoU为0.6的两个框，若得分为0.9与0.7，阈值为0.5，第二个会被删除。这一规则处理重复预测，却可能伤害拥挤场景。一个单元共享类别预测也限制了原始版本处理同单元多对象的能力；后来的YOLO版本改变了预测、分配与去重设计，不能沿用此处全部细节。

\subsubsection{多尺度特征与密集训练权重：FPN、RetinaNet}
\label{sec:vision-recognition-fpn}

远处车辆在图上只占几个像素，深层特征的空间步幅可能使它消失。浅层保留更多位置细节，深层又具有较大的上下文。\term{特征金字塔网络}（feature pyramid network, FPN）把深层特征上采样，与相应浅层经通道变换后的特征相加，再用卷积整理，形成不同空间尺度的预测特征\scite{vision-1612-03144}

第$\ell$层的步幅$s_\ell$连接特征坐标与原图坐标。训练需要按使用该金字塔的检测器规定，将目标或锚框分配到尺度层；FPN本身只生成特征，并不规定一种通用的最终检测集合。较密特征有助于保留小目标，但若目标在输入缩小时已经丢失，特征融合不能补回观测。

RetinaNet在金字塔各位置放置预定大小和长宽比的\term{锚框}（anchor box），即坐标回归的参考框。网络预测类别以及相对于参考框的中心偏移、宽高对数修正。训练按锚框与真实框的IoU分配前景、背景和忽略项，推理则解码修正、筛选并执行NMS。

大量背景锚框在训练后很容易正确分类，但数量仍可能主导总损失。令$p_t$为模型分给真实二值标签的概率，\term{焦点损失}（focal loss）为
\begin{equation}
 L_{\mathrm{focal}}=-\alpha_t(1-p_t)^\gamma\log p_t\text{，}
 \qquad\gamma\geq0\text{。}
 \label{eq:vision-focal-loss}
\end{equation}
其中$\alpha_t$调整正负样本权重。$\gamma=0$时退回加权交叉熵；$\gamma=2$时，$p_t=0.99$的容易样本乘上$10^{-4}$，而$p_t=0.5$的困难样本乘上$0.25$。这个差别让训练更关注尚未学好的候选\scite{vision-1708-02002}
改变容易背景的权重没有增加少见类别的观测，因而前景背景失衡与类别长尾仍须分别检查。

\subsubsection{逐位置边界距离预测：FCOS}
\label{sec:vision-recognition-fcos}

预定义参考框带来了大小与比例的设计。FCOS改为让特征位置直接预测到对象左、上、右、下边界的距离$(l,t,r,b)$。若位置对应原图$(x,y)$，预测框就是$(x-l,y-t,x+r,y+b)$，不再以锚框为坐标起点\scite{vision-1904-01355}

监督仍然需要分配。位于真实框内且满足尺度范围的位置可以成为正例；原方法用不同层的回归范围分摊对象，并在多个框竞争一个位置时选择面积较小的框。负例只学习分类，正例还学习距离及中心度。这里描述的是原始FCOS，后续采用中心区域采样等改动时应另记配置。

\term{中心度}（centerness）衡量位置是否接近框中心，训练目标为
\begin{equation}
 c_*=\sqrt{\frac{\min(l,r)}{\max(l,r)}
                 \frac{\min(t,b)}{\max(t,b)}}\text{。}
 \label{eq:vision-fcos-centerness}
\end{equation}
式中距离来自分配的真实框，框内部距离为正。中心点四侧相对均衡，$c_*=1$；靠近边界时其中一个比例趋于零。推理把预测中心度与类别得分组合，降低边缘候选的排序，再用NMS形成集合。

对于框$(20,10,80,50)$，中心$(50,30)$的目标距离为$(30,20,30,20)$。位置$(25,30)$的距离变成$(5,20,55,20)$，中心度为$\sqrt{1/11}\approx0.302$。即使这两个位置都能精确回归同一框，中心度也不同，说明它不是预测框IoU的直接测量。图\ref{fig:vision-box-points}把位置、距离与框的关系画在同一坐标系中。

\begin{figure}[htbp]
 \centering
 % 连续像素坐标下的FCOS教学例；所有长度与正文算例一致。
\begin{tikzpicture}[x=1mm,y=-1mm,font=\figurefont\small,
  dist/.style={-{Stealth[length=2mm]},draw=BookBlue,line width=0.9pt}]
 \draw[draw=BookMuted,line width=0.7pt,->] (10,0)--(90,0) node[right] {$x$};
 \draw[draw=BookMuted,line width=0.7pt,->] (10,0)--(10,58) node[below] {$y$};
 \draw[draw=BookTeal,fill=BookTeal!7,line width=1pt] (20,10) rectangle (80,50);
 \draw[draw=BookMuted,dashed,line width=0.7pt] (50,10)--(50,50);
 \draw[draw=BookMuted,dashed,line width=0.7pt] (20,30)--(80,30);
 \fill[BookInk] (50,30) circle (1mm);
 \fill[BookBlue] (25,30) circle (1mm);
 \node[anchor=north west] at (51,33) {中心：$c_*=1$};
 \node[anchor=south west] at (26,27) {边缘：$c_*\approx0.302$};
 \draw[dist] (25,18)--(20,18);
 \draw[dist] (25,18)--(80,18);
 \node[anchor=south] at (22.5,17) {$5$};
 \node[anchor=south] at (57,17) {$55$};
 \draw[draw=BookBlue,line width=0.7pt] (25,18)--(25,29);
 \node[anchor=north] at (20,51) {$(20,50)$};
 \node[anchor=south] at (80,9) {$(80,10)$};
\end{tikzpicture}

 \caption{同一真实框内的中心位置与边缘位置。实线箭头标出边缘位置到左右边界的距离，虚线标出中心轴；两点都可以回归同一矩形，但中心度目标分别为1与约0.302。坐标采用向右、向下为正的连续像素边界。}
 \label{fig:vision-box-points}
\end{figure}

\subsubsection{无序对象集合与唯一匹配：DETR}
\label{sec:vision-recognition-detr}

密集候选需要在输出端去重。DETR把目标直接视为无序集合：固定$n$个对象查询读取图像特征，每个查询预测一个框和类别，类别还包含“无对象”。查询是学习得到的向量，并非预先指定给某辆车的身份\scite{vision-2005-12872}

训练标注有$m\leq n$个对象，记第$i$个真实类别与框为$(c_i,\bm{b}_i)$，第$j$个预测为$(\hat{\bm{p}}_j,\hat{\bm{b}}_j)$。寻找从真实对象到不同预测位置的单射$\sigma$：
\begin{equation}
 \hat{\sigma}=\argmin_\sigma
 \sum_{i=1}^{m}\left[
 -\hat p_{\sigma(i)}(c_i)
 +\lambda_1\lVert\bm{b}_i-\hat{\bm{b}}_{\sigma(i)}\rVert_1
 +\lambda_{\mathrm{g}}L_{\mathrm{GIoU}}
    (\bm{b}_i,\hat{\bm{b}}_{\sigma(i)})\right]\text{。}
 \label{eq:vision-detr-matching}
\end{equation}
框在共同归一化坐标中比较。广义IoU还惩罚两个框最小包围矩形中未被并集覆盖的区域，使无交集时也有几何信号。匹配代价用于选对应；选定后，分类损失采用负对数概率，定位项只计算匹配对象，未匹配查询学习“无对象”，并降低该类的权重。

假设两个真实对象对三个查询的总代价矩阵为
\[
 \bm{C}=\begin{pmatrix}1&2&9\\1.2&5&8\end{pmatrix}\text{。}
\]
逐行独立取最小都会选第一个查询，违反唯一对应。集合匹配选择第一对象到第二查询、第二对象到第一查询，总代价为$3.2$，小于反过来分配的$6$。三辆车的情形同样要求三个不同查询；改变标注排列只改变匹配的表示，不应改变最佳集合代价。

使用时没有真实框参与匹配，直接过滤“无对象”和低分输出。原始DETR不执行常规NMS，重复预测主要由唯一匹配训练约束。查询数限制一次可输出的对象数，原始全局注意力、训练周期与小目标效果也各有条件；这些限制不能概括后来所有Transformer检测器。

\Needspace{65mm}
\subsubsection{文字条件下的对象预测：Grounding DINO}
\label{sec:vision-recognition-grounding-dino}

当目标变成“红色货车”时，需要把类别与属性文字接入检测计算。Grounding DINO分别编码图像和文字，通过跨模态特征增强交换信息，再用语言相关性选择用于初始化的对象查询；解码器继续读取两侧特征并更新框与文字对应得分\scite{vision-2303-05499}

文字经分词形成若干词元，候选图像特征与词元比较，训练用对象范围以及对象对应的词语监督匹配。推理时不仅要还原框，还要说明得分对应“货车”、包含属性的短语，还是更长表达中的一部分。词元阈值、框阈值和短语聚合规则都会影响最终输出。

“红色货车”要求属性绑定，“公交车旁的货车”又要求利用参照物消歧。能接收这两种字符串只说明接口容纳了条件，不能证明关系已经正确解析。可分别遮去参照物、更换颜色词或保留多个同类目标，检查预测是否按所需证据改变。完整的表达定位及证据协议在\ref{chap:vision-multimodal-understanding}章讨论。

\subsection{像素归属与对象区分}
\label{sec:vision-recognition-segmentation}

框内常含背景，也可能含另一辆车。需要准确范围时，应进一步指定像素归属。\term{语义分割}（semantic segmentation）给像素分配类别，同类车辆可以连成一片；\term{实例分割}（instance segmentation）给不同车辆各自的掩码；\term{全景分割}（panoptic segmentation）还共同覆盖道路、天空等非实例区域，使每个像素获得类别，可数对象的像素再获得实例编号。

这一区别决定结果的组织方式。两辆车的并集可以是正确的车辆语义区域，却是错误的单个实例。道路通常按区域类别处理，不需要把每块路面任意编号。检测、语义与实例标签不能未经说明相互转换。

图\ref{fig:vision-segmentation-outputs}沿用同一个离散街景，显示增加实例编号与背景类别各带来什么信息。这里仅画出可见区域，框状图元表示教学掩码，不是算法检测框。

\begin{figure}[htbp]
 \centering
 % 同一离散场景的三种分割输出；字符在灰度下保留类别和身份区别。
\begin{tikzpicture}[font=\figurefont\small,x=4.5mm,y=-4.5mm]
 \foreach \offset/\title in {0/语义分割,8/实例分割,16/全景分割}{
  \begin{scope}[xshift=\offset*4.5mm]
   \node[anchor=south] at (3,0) {\title};
   \fill[BookPaper] (0,0.4) rectangle (6,4.4);
   \draw[BookMuted,line width=0.7pt] (0,0.4) rectangle (6,4.4);
   \draw[BookRule,line width=0.5pt] (0,1.4)--(6,1.4);
   \ifnum\offset=8
    \node[text=BookMuted] at (3,3.6) {不标背景};
   \else
    \node[text=BookMuted] at (3,3.6) {道路};
    \node[text=BookMuted] at (3,0.9) {天空};
   \fi
   \filldraw[fill=BookTeal!18,draw=BookTeal,line width=0.8pt]
      (0.5,1.8) rectangle (2.8,2.8);
   \ifnum\offset=0
    \filldraw[fill=BookTeal!18,draw=BookTeal,line width=0.8pt]
      (3.2,1.8) rectangle (5.5,2.8);
    \node at (1.65,2.3) {车};
    \node at (4.35,2.3) {车};
   \else
    \filldraw[fill=BookAmber!16,draw=BookAmber,dashed,line width=0.8pt]
      (3.2,1.8) rectangle (5.5,2.8);
    \node at (1.65,2.3) {车1};
    \node at (4.35,2.3) {车2};
   \fi
  \end{scope}
 }
\end{tikzpicture}

 \caption{三种分割结果的区别。语义分割保留类别，实例分割分开可数对象，全景分割共同保留背景类别与对象编号；文字编号和边界线型使身份区别在灰度打印中仍可辨认。图中区域为离散教学构造。}
 \label{fig:vision-segmentation-outputs}
\end{figure}

\subsubsection{从空间特征形成像素判断}
\label{sec:vision-recognition-dense-pixels}

密集特征仍比原图稀疏。把它们简单放大只能增加采样位置，细边界需要由保留位置的信息支持。单幅语义图和双时相变化图都以像素为输出，但前者判断类别，后者判断同一区域的状态是否改变。

\Needspace{60mm}
\paragraph{多尺度上下文与边界恢复：DeepLabv3+}
\label{sec:vision-recognition-deeplab}

若减小空间步幅，计算成本随特征尺寸上升；若持续下采样，路杆和车窗边缘可能难以恢复。DeepLabv3+用空洞卷积扩大特征读取范围，结合不同采样间隔的上下文，再以解码器融合较浅层细节\scite{vision-1802-02611}

空洞率$r$表示卷积核相邻采样位置间隔为$r$，例如一维三点核读取$x_{i-r},x_i,x_{i+r}$，参数个数仍为三。多尺度分支与全局汇总使位置可以同时利用邻域和较大范围。解码阶段上采样语义特征，对浅层特征先调整通道，再拼接和卷积，最后恢复输入网格上的类别得分。

训练对有标注像素计算交叉熵，忽略无效标签；推理逐位置选择类别并按预处理的逆坐标变换恢复原图。道路内部的类别主要依赖上下文，细路杆又需要足够分辨率。改变输出步幅和输入尺寸会改变这两类证据与预算，应共同记录。即使两车边缘都清楚，纯类别图仍未分配实例身份。

\paragraph{双时相上下文与变化图：BIT}
\label{sec:vision-recognition-bit}

遥感影像的目标可能是发现新增建筑，而不是给单幅影像中的每栋楼分类。BIT输入同一区域的两期图像，以共享的特征提取器形成空间表示，再把每期特征通过学习的空间权重汇总为少量语义词元。两期词元共同编码，使一幅图的局部判断能参考另一时期的上下文\scite{vision-2103-00208}

可将词元汇总理解为$\bm{T}=\bm{A}^{\top}\bm{F}$：$\bm{F}\in\mathbb{R}^{n\times d}$是$n$个位置的特征，$\bm{A}\in\mathbb{R}^{n\times k}$的每列在位置上归一化，得到$k$个加权汇总。编码后的词元再通过交叉注意力返回两期空间位置，恢复上下文增强的特征；对两期特征作差异计算并用预测头输出变化概率。

训练使用配准影像与像素变化标签，模型参数由变化损失更新。新建筑应对应语义变化，树木变绿是否计为变化则取决于标注定义。若两期建筑整体错开几个像素，差异分支也会产生边缘响应，因此配准是任务前提。上下文交换能够利用关联，却不能把传感器差异自动解释成真实建设活动。

\subsubsection{从对象或区域形成掩码}
\label{sec:vision-recognition-mask-sets}

为了分开同类对象，可以先取得候选对象，再在对象内预测掩码；也可以直接预测一组掩码并匹配实例。用户点框提示又提供了另一种目标选择依据。以下沿这三种接口展开。

\paragraph{区域对齐与实例掩码：Mask R-CNN}
\label{sec:vision-recognition-mask-rcnn}

Mask R-CNN在候选框上并行预测类别、框修正和实例掩码。为了让框内位置与原特征准确对应，\term{区域对齐}（RoIAlign）保留连续坐标，以双线性插值在各小格中采样，再汇总成固定尺寸特征，避免先把框边界粗略取整造成的位移\scite{vision-1703-06870}

设四个相邻特征值为$f_{0,0},f_{0,1},f_{1,0},f_{1,1}$，连续位置在水平方向偏移$a$、竖直方向偏移$b$，二者均在$[0,1]$内，则插值为
\begin{equation}
 f(a,b)=(1-a)(1-b)f_{0,0}
 +a(1-b)f_{0,1}+(1-a)bf_{1,0}+abf_{1,1}\text{。}
 \label{eq:vision-roialign}
\end{equation}
用数值$(0,4,2,6)$与$(a,b)=(0.25,0.5)$可得$f=2$，既保留横向的四分之一偏移，也保留纵向中点。若直接取最近格点，这些偏移信息就会消失。

原方法对每个类别预测一个二值掩码通道。训练时只对该正候选匹配到的真实类别通道计算像素二元损失，其他类别通道不在此样本上竞争像素归属。使用时取预测类别的掩码，将区域坐标缩放回原框并放回原图。两个重叠车辆由两个候选各自形成掩码，重叠区域的实例归属与类别判断是不同问题。

区域对齐解决采样量化误差，不能补回候选阶段漏掉的车辆，也不能修正把候选匹配到错误实例的监督。其关键点扩展改变输出目标为有部位语义的位置，不能由掩码准确直接推断关节准确。

\paragraph{掩码查询与局部注意力：Mask2Former}
\label{sec:vision-recognition-mask2former}

先有框再分割并非唯一方式。Mask2Former从多尺度像素特征与查询直接产生掩码集合。每个查询预测类别以及掩码向量，掩码向量与各位置像素嵌入比较，形成该查询的前景概率；训练用集合匹配连接预测与目标，再计算分类、二元掩码和Dice损失\scite{vision-2112-01527}

Dice项度量软预测与真实前景的重叠。对于预测$p_j$和真实标签$m_j$，一种带平滑量$\epsilon>0$的形式为
\begin{equation}
 L_{\mathrm{Dice}}=
 1-\frac{2\sum_j p_jm_j+\epsilon}
 {\sum_j p_j+\sum_j m_j+\epsilon}\text{。}
 \label{eq:vision-dice}
\end{equation}
它与逐像素损失提供不同信号：一个关注区域重叠，另一个约束各位置的预测。训练中的采样点、归一化及匹配权重还需按具体实现固定。

掩码注意力把上一轮预测的前景区域变为后续跨注意力的可读范围。对于查询$i$与位置$j$，在注意力得分中加入屏蔽项$M_{i,j}$：允许读取为0，不允许为$-\infty$，再作softmax。这样每个查询重点读取自己可能对应的区域；多尺度特征与反复更新使区域可以细化。

输出相同的“类别加掩码”接口，可以按类别汇总成语义图，也可保留实例，或通过类别和重叠处理形成全景图。这些恢复规则和训练标签不同，原工作按任务与数据分别训练，不能把统一架构理解为一份参数无条件通用。错误的早期掩码也可能限制后续读取，因此要检查小对象和被遮挡对象是否在迭代中持续被漏掉。

\paragraph{图像与提示编码、歧义掩码：SAM}
\label{sec:vision-recognition-sam}

有时用户只希望选中图上的一个区域，甚至尚未给它类别名。SAM把图像编码、提示编码与掩码解码分开，同一图像的昂贵编码可以复用，点或框变化时更新提示并运行较轻的解码器。点提示包含位置与前景、背景标签，框提示提供空间范围；已有掩码也可作为稠密提示\scite{vision-2304-02643}

解码器让提示与图像表示交换信息，输出若干候选掩码及质量估计。训练提供图像、目标掩码与采样提示，歧义输出通过多个候选容纳一个提示的不同合理解释。质量预测估计候选与监督目标的重叠程度，服务选择，并非在使用时知道了真实掩码。

落在车门上的正点可能指车门或整车。用户选择接近意图的候选，再在漏掉的车尾补正点，或在被误选的路面补负点，才逐步消除歧义。应报告提示来源、点击次数与纠错策略；用真实标注框作提示与由检测器自动给框，具有不同的信息与错误条件。SAM的区域输出本身不自动附带“货车”名称。

\paragraph{概念存在与实例范围：SAM 3}
\label{sec:vision-recognition-sam3}

点提示通常选择一个局部区域，“找出所有校车”则需要判定概念是否出现，并枚举符合概念的实例。SAM 3接收概念短语和带正负标签的图像区域示例，区域视觉特征参与概念表示；随后由对象查询预测框、掩码和局部匹配得分。\nocite{vision-2511-16719}\href{https://arxiv.org/abs/2511.16719}{SAM 3原始报告}将这一接口称为可提示概念分割。

概念是否在整幅图像中出现，与每个候选是否匹配，承担不同职责。共享存在词元与对象查询一同进入解码器，再由存在预测头得到$p_{\mathrm{pres}}$；查询$i$的局部分数$s_{i\mid\mathrm{pres}}$表达概念存在条件下的匹配，最终使用$s_i=p_{\mathrm{pres}}s_{i\mid\mathrm{pres}}$。训练分别提供存在标签与实例范围；局部分类损失在概念存在的图像上计算，使全局缺席主要由共享分支学习。不能仅用“图里有车”这一标签学出所有车辆的像素边界。

对于不存在的概念，即使某个候选与“车”这一宽泛外观相似，共享存在判断也可以抑制错误输出。反过来，共享存在分支若把“校车”判为不存在，多个局部查询都可能被压低，形成相关的漏检。评价因此分别检查概念有无、每实例召回和掩码边界。简短概念与区域示例扩大了目标选择方式，但尚不足以推断任意关系嵌套或长指令都已解析；视频身份传播在\ref{chap:vision-video-recognition}章继续讨论。

\subsection{软边界与前景贡献的估计}
\label{sec:vision-recognition-matting}

发丝和半透明边缘的一个像素可能同时收到前景与背景的光。二值掩码只能选择保留或删除，无法表达这种混合。\term{图像抠图}（image matting）估计前景贡献比例$\alpha\in[0,1]$。在线性颜色空间的简化合成模型中，每个像素满足
\begin{equation}
 \bm{x}=\alpha\bm{f}+(1-\alpha)\bm{b}\text{，}
 \label{eq:vision-matting}
\end{equation}
其中$\bm{x},\bm{f},\bm{b}\in\mathbb{R}^3$是观测、前景与背景颜色。只有三个颜色观测，却有前景、背景和$\alpha$共七个未知标量，需要额外条件或先验；二值分割标签并不能唯一解出它们。

\subsubsection{三分图条件与alpha细化：Deep Image Matting}
\label{sec:vision-recognition-deep-matting}

Deep Image Matting使用\term{三分图}（trimap）指定确定前景、确定背景和未知区域，把RGB图与这一提示通道共同输入编码解码网络，再用细化网络改善$\alpha$边界。确定前景约束$\alpha=1$，确定背景约束$\alpha=0$，学习的重点是未知带中的混合比例\scite{vision-1703-03872}

训练可利用已知前景、背景与$\alpha$构造合成图像，监督$\alpha$误差，并按式\eqref{eq:vision-matting}检查预测贡献比例能否重建合成观测。使用时输出$\alpha$及其已知区约束；要把对象放到新背景$\bm{b}'$，仍需估计或取得前景颜色$\bm{f}$，再合成为$\alpha\bm{f}+(1-\alpha)\bm{b}'$。只取得$\alpha$并不意味着前景颜色已无污染。

例如灰度前景为0.8、原背景为0.2、$\alpha=0.25$，观测为0.35。换成0.6的背景，正确合成为$0.25\times0.8+0.75\times0.6=0.65$；若二值化为背景，结果为0.6，发丝的贡献消失。若直接把观测0.35当成前景颜色，又会混入旧背景。提示误把发丝标成确定背景时，网络还可能被错误条件限制。边界、贡献比例和最终合成质量应分别核验。

\section{数量、姿态与空间测量}
\label{sec:vision-recognition-measurement}

“有几辆车”“人的手腕在哪里”“路牌有多远”分别要求数量、具有部位语义的位置以及带坐标或尺度的空间量。它们可以调用前面的对象范围，但监督和误差不能由检测指标替代。

\subsection{对象数量的估计}
\label{sec:vision-recognition-counting}

稀疏对象可以先逐个识别再计数，密集遮挡场景则可以直接估计数量的空间分布。前者保留对象清单，后者强调总量，两者取得正确总数的原因可能不同。

\subsubsection{检测实例汇总与重复计数控制}
\label{sec:vision-recognition-instance-counting}

给定去重后的检测集合，可按类别、区域与可见条件筛选，再把满足条件的指示量相加。若区域为$R$，以框中心$\bm{c}_i$是否在$R$内为计入规则，则
\begin{equation}
 \hat n_R=\sum_i
 \mathbf{1}\{c_i\in C_{\mathrm{target}}\}
 \mathbf{1}\{s_i\geq\eta\}
 \mathbf{1}\{\bm{c}_i\in R\}\text{。}
 \label{eq:vision-detection-count}
\end{equation}
这里$c_i$表示类别，$\bm{c}_i$表示二维中心，两者字形不同。中心落入与“任意可见部分落入”是两种不同规则，边界目标会得到不同计数。

例如真实有5辆车，漏掉1辆，另1辆被重复预测，总数仍可能为5。只看计数误差会掩盖两种错误；应同时检查对象匹配，尤其是遮挡和边界处的实例。跨帧统计还要避免同一辆车被多次计入，需由跟踪维持身份。

\subsubsection{多尺度密度预测与积分：MCNN}
\label{sec:vision-recognition-mcnn}

在人群过密、单个身体范围难以标注时，可用头部点标注学习\term{密度图}（density map），即每个位置分摊多少计数贡献。MCNN以多个不同感受范围的卷积分支处理尺度变化，连接各列特征后回归一张密度图\scite{vision-mcnn2016}

若标注点为$\bm{p}_i$，可用归一化核构造监督
\begin{equation}
 D(\bm{p})=\sum_{i=1}^{n}G_{\sigma_i}(\bm{p}-\bm{p}_i),
 \qquad
 \sum_{\bm{p}}G_{\sigma_i}(\bm{p}-\bm{p}_i)=1\text{。}
 \label{eq:vision-density-target}
\end{equation}
原工作令核宽与邻近头部点的平均距离相关，用局部拥挤程度近似尺度变化。卷积列和汇总层以预测密度与目标密度的平方误差训练。使用时对预测网格求和得到人数；密度可以是非整数分布，最终总数也可保留实值用于评价。

三个点各贡献1，整图密度之和便为3。若一个边缘点的核有一半落到图外，直接截断后该点只贡献约0.5；有限网格构造标签时需按计入规则归一化有效核。图像缩放也要保持总和：将每个密度单元分成四个，四个新单元应分享原来的质量，不能每个都复制原值。

总和正确仍可能把本属左半边的人数移到右半边。应同时看局部计数、空间误差与场景变化。透视改变对象尺度，新的拍摄高度和密度范围可能超出训练条件；密度图没有自动恢复每个人的身份或完整位置。

\subsection{关节与姿态的估计}
\label{sec:vision-recognition-pose}

人体关节不同于图像中任意可匹配的角点：左腕、右肘具有固定部位含义。多人场景还必须知道这些部位属于谁。二维位置位于图像平面，三维姿态则要增加深度与坐标约定。

\subsubsection{二维部位定位与关联：OpenPose}
\label{sec:vision-recognition-openpose}

自顶向下路线先检测人，再在各人区域内定位关节；多人重叠时，候选框和裁剪质量影响后续姿态。OpenPose采用自下而上路线，在整图预测关节热图以及连接部位的\term{部位亲和场}（part affinity field, PAF），再组装人体\scite{vision-1611-08050}

热图在对应关节附近取得高响应，峰值形成候选位置。PAF是每种肢段的一张二维向量场：真实肢段附近写入从一种部位指向另一种部位的单位向量，其他位置为零，重叠肢段按规则合并。训练从关键点标注构造两类目标，对有效位置计算热图与向量场误差，并逐阶段细化。

设候选肘部、腕部为$\bm{a},\bm{b}$，方向$\bm{u}=(\bm{b}-\bm{a})/\lVert\bm{b}-\bm{a}\rVert_2$，对应PAF为$\bm{L}$。沿候选肢段平均方向一致程度：
\begin{equation}
 s(\bm{a},\bm{b})=
 \int_0^1
 \bm{L}((1-\lambda)\bm{a}+\lambda\bm{b})^\top
 \bm{u}\,\mathrm{d}\lambda\text{。}
 \label{eq:vision-paf-score}
\end{equation}
实现中在有限个点上采样求平均，再按部位连接类型作匹配并组合人体。$\bm{a}=\bm{b}$时方向未定义，应排除这一退化候选。

假设五个采样点的方向内积为$(1,1,0.8,1,0.7)$，平均为0.9；连到另一人的手腕时变为$(0.1,0,-0.2,0,0.1)$，平均为0。两个手腕都定位正确，也仍需这一关联步骤才可能把手臂接对。投影重合、严重遮挡和交叉肢段会使证据混杂，因此关节位置与人体组装应分开评价。

\subsubsection{二维关节到三维姿态：Martinez基线}
\label{sec:vision-recognition-lifting}

取得二维关节后，可以把按固定部位顺序排列的$2j$维坐标输入回归网络，输出$3j$维三维坐标。Martinez等人的基线使用全连接层与残差模块学习这一映射，使二维检测和三维恢复的误差来源能够分开研究\scite{vision-1705-03098}

训练规范化输入坐标，以对应的三维姿态作回归监督。三维目标通常相对于根关节表示，并遵守指定相机坐标、尺度和关节集合。使用时必须复现训练的规范化，再将输出还原为该协议中的坐标；关节顺序错位会使一个部位的输入被当成另一个部位，即使维度仍然匹配。

相同的二维投影可以来自向前或向后弯曲的肢段，网络用训练中的人体结构先验选择一种三维解释。真实二维关节作为输入可以揭示提升网络本身的限制，检测二维关节则包含两阶段误差。评价要说明是否允许旋转、平移及尺度对齐：相似变换对齐后的关节误差较小，不意味着原始输出已具有正确朝向或距离。多视角三角化增加了观测约束，但需要可靠相机和跨视角关节对应。

\subsection{深度与距离的估计}
\label{sec:vision-recognition-depth}

从二维走向三维，需要知道相机如何把空间点投到图像上。在忽略畸变的针孔模型中，世界坐标$\bm{P}$通过旋转$\bm{R}$、平移$\bm{t}$进入相机坐标，再由内参矩阵$\bm{K}$投影：
\begin{equation}
 z\begin{pmatrix}u\\v\\1\end{pmatrix}
 =\bm{K}(\bm{R}\bm{P}+\bm{t})\text{。}
 \label{eq:vision-camera-projection}
\end{equation}
其中$z$为相机光轴方向的深度，$(u,v)$为像素坐标。深度与到相机中心的欧氏距离通常不同；偏离光轴的点，其距离还包含横向分量。

单个像素约束一条射线而非唯一空间点。把所有三维坐标及相机平移同时乘正数$a$，式\eqref{eq:vision-camera-projection}中的齐次比例也乘$a$，像素位置不变。缺少已知基线、物体尺寸或其他度量信息时，这说明尺度不能仅由投影唯一决定。

\subsubsection{单幅相对深度与度量适配：Depth Anything V2}
\label{sec:vision-recognition-depth-anything}

单目模型利用物体形状、透视、遮挡和训练先验，从一张图预测密集深度。Depth Anything V2以合成数据的精确深度训练较强教师，再为大量真实图像产生深度伪标签，供不同规模学生学习；图像编码器与稠密解码器把表示接回各空间位置\scite{vision-2406-09414}

这一训练桥接应区分两个信息来源：合成图像提供几何标签，真实图像提供目标外观，伪标签质量依赖教师。相对预测可通过尺度、平移不变的目标学习排序和形状；实际输出可能采用深度或逆深度相关约定，使用时应核查具体检查点，而不能仅按图像颜色解释米数。

若三个位置的相对逆深度为$(0.8,0.4,0.2)$，可以按该约定判断第一个较近，却不能未经校准说深度为$(1.25,2.5,5)$米。用带实际单位的深度标签适配后，模型才学习特定数据条件下的度量映射；即便如此，相机和场景变化仍可能带来尺度误差。边界位置、前后排序和实际距离应分别评价。

\subsubsection{相机条件下的代价体与深度推断：MVSNet}
\label{sec:vision-recognition-mvsnet}

如果有多张重叠照片和已知相机，同一表面在不同图像中的对应可以提供深度约束。MVSNet选定参考图，对每个候选深度$d_k$，把参考像素反投影成空间点，再投到其他相机，在源图特征中插值取值。整个平面上的映射可写为依赖$d_k$的可微单应变换\scite{vision-1804-02505}

正确深度应使同一表面的特征比较一致。将参考特征和变换后的各源特征记为$\bm{f}_r$，共有$n_{\mathrm{v}}$个视图，可以用方差形成匹配代价：
\begin{equation}
 \bar{\bm{f}}=\frac1{n_{\mathrm{v}}}\sum_{r=1}^{n_{\mathrm{v}}}\bm{f}_r,\qquad
 \bm{c}(d_k,u,v)=\frac1{n_{\mathrm{v}}}\sum_{r=1}^{n_{\mathrm{v}}}
   (\bm{f}_r-\bar{\bm{f}})\odot(\bm{f}_r-\bar{\bm{f}})\text{。}
 \label{eq:vision-mvs-cost}
\end{equation}
这里$\odot$表示逐元素相乘。所有像素和候选深度的代价共同形成\term{代价体}（cost volume），它保留“每个位置在各个深度假设下多不一致”的信息。三维卷积在空间与深度轴上整理代价，经softmax得到深度权重$\pi_k$，用$\hat d=\sum_k\pi_kd_k$回归深度，并可作后续细化。训练只对具有有效参考深度的位置计算误差。

在平行、已校正的双目简化情形中，焦距$f$、相机基线$b$和水平视差$\Delta u$满足$d=fb/\Delta u$。取$f=800$像素、$b=0.2$米、$\Delta u=40$像素，得到$d=4$米。视差误差1像素在此处约引起0.1米深度误差；同样误差在更远处会被放大。这一关系说明已知基线如何提供实际尺度。

图\ref{fig:vision-depth-hypotheses}展示两个深度假设如何查询不同源图位置。弱纹理可能让多个候选都具有相近代价，遮挡使源图看到另一表面，反光破坏外观一致性，相机误差又使正确深度也无法对齐。需要同时报告深度误差、可恢复面积与候选数带来的存储成本。

\begin{figure}[htbp]
 \centering
 % 针孔投影的几何关系；示意深度候选，不按物理长度绘制。
\begin{tikzpicture}[font=\figurefont\small,x=1cm,y=1cm,
 ray/.style={draw=BookBlue,line width=0.9pt},
 alt/.style={draw=BookAmber,dashed,line width=0.9pt}]
 \coordinate (a) at (0,0);
 \coordinate (b) at (7,0);
 \coordinate (p) at (2,3);
 \coordinate (q) at (3,4.5);
 \draw[ray] (a)--(q);
 \draw[ray] (b)--(p);
 \draw[alt] (b)--(q);
 \draw[draw=BookInk,line width=1.2pt] (-0.3,1)--(2.2,1);
 \draw[draw=BookInk,line width=1.2pt] (4.5,1)--(7.5,1);
 \fill[BookInk] (a) circle(1.1mm);
 \fill[BookInk] (b) circle(1.1mm);
 \fill[BookTeal] (p) circle(1.1mm);
 \draw[BookAmber,fill=white,line width=0.9pt] (q) circle(1.1mm);
 \fill[BookBlue] (0.667,1) circle(1mm);
 \fill[BookBlue] (5.333,1) circle(1mm);
 \draw[BookAmber,fill=white,line width=0.9pt] (6.111,1) circle(1mm);
 \node[below] at (a) {参考相机};
 \node[below] at (b) {源相机};
 \node[above right] at (0.667,1) {同一参考像素};
 \node[below] at (5.333,0.9) {$u'_1$};
 \node[below] at (6.111,0.9) {$u'_2$};
 \node[left] at (p) {深度$d_1$};
 \node[right] at (q) {深度$d_2$};
 \node[right] at (7.5,1) {源图平面};
 \draw[draw=BookMuted,<->,line width=0.7pt] (0,-0.8)--(7,-0.8)
   node[midway,below] {已知相机相对位置};
\end{tikzpicture}

 \caption{已知相机下，参考像素沿射线形成不同深度假设，分别投到源图的不同位置。图为针孔几何教学示意，只表示对应关系，不按物理长度绘制。跨图比较可排除不一致的假设，但遮挡或重复纹理会使判断不唯一。}
 \label{fig:vision-depth-hypotheses}
\end{figure}

\subsection{相机与场景结构的恢复}
\label{sec:vision-recognition-geometry}

前一方法把相机作为已知条件。若相机未知，需要共同解释视图关系、表面点和投影。经典流程先寻找对应，再估计相机和三维点；学习式方法可以直接预测这些几何量，仍需处理共同坐标与尺度。

\Needspace{65mm}
\subsubsection{成对点图与全局对齐：DUSt3R}
\label{sec:vision-recognition-dust3r}

DUSt3R输入图像对，以两侧编码及跨图解码交换信息，为两幅图的像素分别输出三维\term{点图}（point map）：每个有效像素对应一个三维向量。两张点图都用第一幅相机的坐标表达，并预测每点置信度，使无需预先给出相机的成对几何成为直接输出\scite{vision-2312-14132}

训练把预测与真实点图按规定尺度规范化，约束坐标误差，并使置信度参与权衡。其基本思想可写为$\sum_i c_i e_i-\alpha\log c_i$，其中$e_i$是点误差，$c_i>0$是预测置信度，$\alpha>0$控制惩罚。只乘$c_i$会使全部置信度趋于零；对数项阻止通过无条件放弃全部点来降低目标。这里的置信度是学习的可靠性权重，不能当作经过概率校准的正确率。

三张建筑照片形成多对输入后，每一对以自己的首图为参考，不能直接把全部点拼成场景。全局对齐优化各对点图到共同坐标的旋转、平移和尺度，使共享视图及可信点彼此相容。相机、深度与对应再从已对齐的几何联系恢复；一次成对前馈与多图优化是两个不同计算阶段。

若甲乙图预测建筑宽度为2个单位，乙丙图为5个单位，需要先估计约2.5的相对尺度并对齐朝向位置。单对看似合理不保证多对闭环一致；弱重叠、动态行人和错误高置信点都会干扰对齐。没有额外度量依据时，共同场景仍只具有约定尺度。

\subsubsection{多视图共享特征与联合几何输出：VGGT}
\label{sec:vision-recognition-vggt}

逐对预测再优化可能重复处理同一视图。VGGT让多幅图的词元共同参与计算，交替进行帧内注意力与跨全部视图的全局注意力。帧内计算保留各图的结构，全局计算交换视角证据，随后由不同预测头输出相机、深度和点图；点跟踪头还接收参考图中的二维查询点\scite{vision-2503-11651}

训练使用对应几何量监督各预测头，把相关任务接在共享表示上。输出以首相机为参考，并沿用训练的尺度规范；这使多视图结果可在同一约定下比较，却没有自动赋予米单位。相机头给出的内外参和深度头的输出可以反投影成点，再与直接点图比较，从而检查联合输出是否几何一致。

给定一个参考点，跨图跟踪输出位置与可见性；没有查询点时，不能把该接口理解为自动枚举全部对象轨迹。多图前馈输出之后仍可进行几何优化，但应把这部分额外时间单独记录。图像数和每图词元数共同增加全局注意力成本，动态对象、超出训练范围的相机以及缺少重叠的图像也会削弱预测依据。新视角渲染将在\ref{sec:vision-generation-novel-view}节使用这些几何条件，生成外观与恢复可测量表面仍是不同目标。

\section{图像匹配与检索}
\label{sec:vision-recognition-retrieval}

识别已经把图像变成类别、范围和测量。检索还需要依据查询，从图库找出相关图像并排序。紧凑描述适合筛选候选，局部对应保留细节，几何约束再检查对应是否能共同解释视角变化。

\subsection{查询与相关性的确定}
\label{sec:vision-recognition-relevance}

“找类似的红色汽车”和“找同一辆车的其他照片”使用不同正例。前者容许不同实例，只要外观或语义相近；后者需要车身细节、车牌等区分身份的证据，并可能允许颜色因光照变化。若没有预先规定相关性，高相似度排序甚至可能朝错误目标优化。

查询可以是整图或裁剪区域。整幅照片包含建筑、树木和天空，局部查询只指定一块招牌；两者对应的容许背景、遮挡和视角变化不同。图库中的近重复图也可能使任务过于容易。评价前应固定查询范围、正例定义、忽略项、图库版本和干扰图规模，文字查询则由\ref{chap:vision-multimodal-understanding}章继续连接。

\subsection{候选图像的形成}
\label{sec:vision-recognition-candidates}

在大图库中逐一做精细匹配成本很高。全局向量先快速保留较少候选，局部特征再进入后续核验。近似索引可以加速向量搜索，但它引入的召回损失与表示本身的区分能力应分开检查。

\subsubsection{全局特征聚合与检索学习：GeM}
\label{sec:vision-recognition-gem}

空间特征需要汇总为固定维向量。平均池化平等计入各位置，最大池化只保留最强响应。\term{广义均值池化}（generalized-mean pooling, GeM）在两者之间调节：对第$k$通道的非负特征$x_{p,k}$，计算
\begin{equation}
 g_k=\left(\frac1n\sum_{p=1}^{n}x_{p,k}^{\,q_k}\right)^{1/q_k},
 \qquad q_k>0\text{。}
 \label{eq:vision-gem}
\end{equation}
指数$q_k=1$为平均，趋于无穷时趋近最大值。实现通常在幂运算前给特征加下界，避免零点数值问题；指数可以固定，也可与编码器参数一同学习\scite{vision-1711-02512}

对响应$(1,1,1,4)$，平均为1.75，$q=3$时为$(67/4)^{1/3}\approx2.559$，最大值为4。提高指数让局部强响应更影响全图描述，但并不保证该响应一定属于目标。聚合后将$\bm{g}$归一化为单位向量，以欧氏距离或内积排序；单位向量满足$\lVert\bm{u}-\bm{v}\rVert_2^2=2-2\bm{u}^{\top}\bm{v}$，两种排序因此等价。

检索训练需要知道哪些图拍到了同一实例。原工作利用三维重建与相机几何选择具有重叠内容的正例和困难负例，故“不需人工检索标注”并不等于没有监督依据。对比目标使正图对接近，对距离小于间隔$m$的负图对施加惩罚；训练后保留编码器与池化参数，为图库和查询使用相同的特征处理。

同一建筑的两张照片可能因视角差异较远，另一栋同样玻璃幕墙的建筑却可能很近。全局描述提供候选，不能证明局部细节一致。已有DINOv3向量可以作为不同参数来源的筛选基线，但比较应固定输入、归一化、图库与训练信息。

\Needspace{55mm}
\subsubsection{局部关键点与描述：SuperPoint}
\label{sec:vision-recognition-superpoint}

若希望保留窗角或招牌角等可重复的位置，可以为每幅图输出稀疏点坐标和描述向量。SuperPoint使用共享编码器连接兴趣点分支与描述分支；兴趣点没有“左腕”那样的对象部位语义，而是适合跨图重复辨认的位置\scite{vision-1712-07629}

原始检测头在较粗网格上预测每个$8\times8$区域内的64个位置及一个“没有点”的通道。经过通道归一化、去掉无点通道和空间重排，得到原图分辨率的点分数，再阈值筛选并抑制过近峰值。描述头产生较稠密的低分辨率特征，在保留点位置插值采样，并将描述归一化。最终保存$(x,y,\text{得分},\bm{d})$，而不仅是一张响应图。

训练先以具有已知角点的合成几何图形学习检测，再对真实图像施加多次已知单应变换，将检测结果逆变换并汇总，形成更稳定的伪标签。已知变换也提供两图描述应对应的位置，用于联合训练点与描述。单应适应提供的是这些变换下的监督，不能覆盖任意三维视差。

窗格重复时，检测器可以在许多窗角上都稳定响应，但描述仍可能相似。因此“点能重复检测”与“点能匹配到正确窗格”是两个要求。尺度、旋转、遮挡及保留点数都会影响后续候选，比较匹配器时需固定这些输入条件。

\subsection{精细匹配与结果排序}
\label{sec:vision-recognition-verification}

候选图像中可能混入外观相似的其他实例。局部描述取得候选点对，上下文比较帮助消歧，几何一致性再排除不能共同成立的对应。几何通过也仍需回到查询定义判断相关性。

\Needspace{55mm}
\subsubsection{上下文局部配对与自适应计算：LightGlue}
\label{sec:vision-recognition-lightglue}

独立最近邻匹配只比较两个描述，无法利用邻近窗格怎样共同移动的信息。LightGlue接收已有关键点和描述，交替通过图内自注意力与跨图信息交换更新表示，再预测匹配关系；它不负责从像素检测关键点\scite{vision-2306-13643}

设两图分别有$n,m$个点，将更新后的描述投影为$d$维，再以$1/\sqrt d$缩放两两内积，得到$S_{i,j}$。两侧点可找到对应的分数为$\mu_i,\nu_j\in[0,1]$，匹配得分结合双向归一化与可匹配性：
\begin{equation}
 P_{i,j}=\mu_i\nu_j
 \frac{\exp S_{i,j}}{\sum_{j'=1}^{m}\exp S_{i,j'}}
 \frac{\exp S_{i,j}}{\sum_{i'=1}^{n}\exp S_{i',j}}\text{。}
 \label{eq:vision-lightglue-score}
\end{equation}
第一项归一化要求$i$在对面候选中偏好$j$，第二项要求$j$在本侧候选中偏好$i$。保留互为最高分且超过阈值的点对，有助于避免多个点争抢同一对应。训练使用真实对应监督匹配得分，对没有对应的点监督可匹配性。

若$\mu_i=0.9,\nu_j=0.8$，两个方向的归一化分数为0.7和0.6，则$P_{i,j}=0.3024$。只有单方向相似度高，或其中一点已被遮挡时，最终得分都可能降低。这是模型的配对评分，不是对几何正确性的独立证明。

不同图像对的难度差异很大。LightGlue在各层预测当前配对是否已稳定，足够比例的点达到置信要求时提前退出；对于已经高置信、但可匹配分数低的点，停止后续计算。尚不能判断的点仍应保留，不能把“不确定”直接当成“无对应”。这些门控调节处理层数和保留点数，比较速度时需固定特征、图像尺寸、硬件与难度分布。

\subsubsection{鲁棒几何估计与候选重排：RANSAC路线}
\label{sec:vision-recognition-ransac}

重复纹理可能产生看似可信但错误的点对，需要检查一组对应能否共同解释视角变化。平面上的点在两图间可由\term{单应变换}（homography）联系，齐次坐标满足$\tilde{\bm{p}}'\sim\bm{H}\tilde{\bm{p}}$。一般三维场景的两视图则遵守\term{极线约束}（epipolar constraint），即$\tilde{\bm{p}}'^\top\bm{F}\tilde{\bm{p}}=0$，其中$\bm{F}$为基础矩阵。后者约束另一图中的一条线，而非单个位置。

\term{随机采样一致性}（random sample consensus, RANSAC）通过反复取少量对应估计模型，寻找被许多点共同支持的解。以单应为例，四个处于非退化配置的点对可以产生候选$\bm{H}$；再把所有对应按重投影残差与阈值分为内点和外点，保存支持较好的模型，并利用内点重估。

若每次独立抽取$s$个对应，内点比例近似为$w$，全内点样本概率约为$w^s$。为了至少以概率$p$抽到一次全内点样本，次数满足
\begin{equation}
 k\geq\frac{\log(1-p)}{\log(1-w^s)}\text{。}
 \label{eq:vision-ransac-iterations}
\end{equation}
这是独立、近似固定内点率下的采样计算，还没有保证样本不退化，也没有保证估计模型正确。取$w=0.5,s=4,p=0.99$，需要至少72次；内点率降低时，所需次数迅速增加。

Graph-Cut RANSAC在局部优化中进一步利用邻近对应的联系。为每个对应设置内外点标签，一元代价来自对当前模型的残差，邻接项倾向于让邻近对应具有一致标签；图割在规定的能量形式下求标记，再据选出的内点重估几何\scite{vision-1706-00984}
图割解决的是给定模型附近的一次局部标记问题，不是对所有模型、样本和标签组合取得全局最优。邻近错误匹配也可能具有一致外观，仍须检查模型假设。

建筑正面近似平面时，单应适用；有明显前后视差时，同一个单应不能同时对齐所有表面，应考虑两视图几何或分区结构。纯旋转、近共线点和重叠太少也会改变估计条件。图\ref{fig:vision-retrieval-evidence}将全局筛选、局部配对与几何核验分别列出，避免把一个阶段的高分误作整条证据链完成。

\begin{figure}[htbp]
 \centering
 % 流程节点间箭头均连接左右中心；下方为各阶段尚未排除的错误。
\begin{tikzpicture}[font=\figurefont\small,
 block/.style={draw=BookTeal,fill=BookTeal!8,minimum width=28mm,
   minimum height=12mm,align=center,inner sep=1.5mm},
 flow/.style={-{Stealth[length=2mm]},draw=BookInk,line width=0.9pt},
 note/.style={align=center,text width=28mm,font=\figurefont\footnotesize,
   text=BookMuted}]
 \node[block] (global) at (0,0) {全局描述\\筛选候选};
 \node[block] (local) at (3.7,0) {局部对应\\匹配细节};
 \node[block] (geo) at (7.4,0) {几何支持\\核验重排};
 \draw[flow] (global.east)--(local.west);
 \draw[flow] (local.east)--(geo.west);
 \node[note] at (0,-1.6) {可能只同类\\并非同一实例};
 \node[note] at (3.7,-1.6) {重复窗格\\可能相互错配};
 \node[note] at (7.4,-1.6) {几何匹配仍需\\满足查询定义};
 \foreach \x in {0,3.7,7.4}
   \draw[draw=BookRule,line width=0.7pt] (\x,-0.65)--(\x,-1);
\end{tikzpicture}

 \caption{图像检索中逐步保留更具体的证据。上方沿主流程减少候选，下方说明每步仍可能留下的错误。几何一致可以帮助重排，却仍要检查所得实例是否满足查询的相关性定义。}
 \label{fig:vision-retrieval-evidence}
\end{figure}

Revisited Oxford和Paris修订了实例检索标注与难度协议，并引入大规模干扰图，适合比较候选筛选和几何重排分别改变了哪些结果\scite{vision-1803-11285}
不同协议对困难正例与忽略项的处理不一样，应固定协议后报告排序。更多内点可以支持同一建筑的判断，却不能直接证明其中招牌文字相同，或建筑具有相同的实际尺度。

\section{结果检验与研究方向}
\label{sec:vision-recognition-evaluation}

组合系统需要让类别、框、掩码、属性和测量指向同一对象。若车辆检测正确，属性却来自旁车，逐模块的总体准确率可能仍然较高。街景应逐实例保留标识并核验空间对应；工业零件应同时检查异常图、整图得分与规定缺陷，避免正常材质变化被业务阈值放大。

检测评价先按类别与IoU阈值对预测和真实对象作一对一匹配，未匹配预测为假阳性，未匹配目标为漏检。按得分变化形成精确率、召回率曲线，再计算平均精度；不同数据集的插值、IoU阈值、尺寸分组与每图检测数量上限必须说明，不能把任意一个“AP”视为同一指标。

区域任务也有不同分母。语义分割按类别比较像素IoU，常再对类别平均；实例分割要先匹配对象。全景质量把匹配区域的IoU与对象错误合在一起，标准任务对同类、IoU超过0.5的区域建立匹配集合$T$，并令$F_{\mathrm{p}},F_{\mathrm{n}}$为多报、漏报区域集合：
\begin{equation}
 \mathrm{PQ}=
 \frac{\sum_{(p,g)\in T}\operatorname{IoU}(p,g)}
 {|T|+\frac12|F_{\mathrm{p}}|+\frac12|F_{\mathrm{n}}|}\text{。}
 \label{eq:vision-pq}
\end{equation}
它同时受到范围与实例错误影响\scite{vision-1801-00868}
例如匹配区域IoU为0.8和0.6，另有一个漏报且无多报，则PQ为$1.4/2.5=0.56$；只报告匹配区域平均IoU会得到0.7，掩盖漏掉的区域。空类别与跨类别平均按数据集协议处理。

检索则按相关性检查排名。若有$R$个相关项，排名$k$处的相关指示为$r_k$，截至该位置的精确率为$P(k)$，全排名上的一种非插值平均精度为
\begin{equation}
 \mathrm{AP}=\frac1R\sum_k P(k)r_k\text{。}
 \label{eq:vision-retrieval-ap}
\end{equation}
三个正例出现在第1、3、5位时，AP为$(1+2/3+3/5)/3\approx0.756$。若只返回前四项，计算截断指标时仍需明示分母和遗漏规则。召回率关注是否找到，AP还关注出现得早不早；无相关项的查询应按协议排除或单独评价，不能直接除以零。

\begin{table}[htbp]
 \centering
 \small
 \begin{tabularx}{\linewidth}{@{}lXX@{}}
 \toprule
 输出 & 需要固定的条件 & 不能由它单独证明的内容\\
 \midrule
 类别、属性 & 类别表、未知项、对象范围 & 对象存在及属性归属都正确\\
 框、掩码 & 匹配规则、提示、可见范围 & 未观测部分及三维形状正确\\
 数量、姿态 & 计入规则、关节与对齐方式 & 实例清单、身份及原始朝向正确\\
 深度、点图 & 尺度、坐标、有效区域 & 距离带实际单位且适合测量\\
 检索排序 & 正例、图库、忽略项及难度 & 对象属性、文字和距离正确\\
 \bottomrule
 \end{tabularx}
 \caption{不同视觉成果所需的评价条件。右列列出仍需另行核验的推论，不表示这些任务必然无法提供相应能力。}
 \label{tab:vision-recognition-protocols}
\end{table}

\Needspace{95mm}
表\ref{tab:vision-recognition-protocols}也说明了研究方向的来源。开放概念增加目标歧义，长尾类别需要更细的覆盖分析，精细边界要求保存小尺度证据，遮挡和单目投影又限制可恢复的空间量。LVIS把大词表和少见实例纳入分割评价，MVTec AD把正常训练与测试缺陷定位分开，为这些问题提供了不同的观察条件\scite{vision-1908-03195,vision-mvtec2019}
医学和遥感等专业用途还要固定采集设备、地理区域或患者划分，核验外观先验能否迁移。模型在一个基准上提供的证据，应停留在该任务与划分实际检验的范围内。

\section{本章小结}
\label{sec:vision-recognition-summary}

从街景得到可用结果，需要把预测接回明确的对象和观测。类别与属性规定语义，框和掩码规定范围，数量与姿态规定测量对象，相机与尺度规定三维解释，检索相关性则决定哪些对应值得保留。共享表示可以支持这些输出，监督、恢复步骤和核验关系仍各自不可省略。

本章的方法提供了两类互补依据。学习表示利用样本规律形成候选判断，空间、几何与集合约束检查这些判断能否共同成立。候选高分不等于目标存在，区域正确不等于身份正确，相对深度不等于米制距离，几何支持也不等于任意业务相关性。只有把这些条件保留到最终输出，才能知道一项结果已经回答了什么，以及仍需要哪些观测。
